\section{L'accréditation des ingénieurs}\label{laccreditation-des-ingenieurs}

\stanceinfo{\textbf{Congrès de la FCEG 2018} [2018-01-07]}{2024-01-02}

\officialstance{La Fédération canadienne étudiante de génie estime que le système d'accréditation des programmes d'ingénierie canadiens doit protéger les intérêt des étudiants en ingénierie en garantissant des normes élevées, cohérentes et réalisables pour la qualité de la formation en ingénierie, en faisant participer les étudiants au processus des visites d'accréditation et à l'élaboration des critères d'accréditation.}

\stanceheading{La position des étudiants}
\begin{itemize}
 \item Les étudiants sont une partie prenante essentielle de la formation des ingénieurs et devraient être activement consultés dans les discussions et les décisions relatives à l'accréditation des ingénieurs.
 \item Le processus de visite d'accréditation devrait être réformé pour mieux évaluer et intégrer les perspectives des étudiants tout au long de la visite.
 \item Un système d'accréditation efficace devrait être principalement axé sur les résultats de l'apprentissage, et toute mesure de la contribution des étudiants devrait être représentative du temps total d'apprentissage et ne pas favoriser l'apprentissage magistral par rapport à d'autres méthodes d'apprentissage.
 \item Les unités d'accréditation (UA) ne favorisent pas toujours les étudiants en ingénierie et n'affectent pas non plus la qualité de la formation d'ingénieur. Il conviendrait d'adopter une approche plus holistique, soit par le biais des UA, soit par une mesure différente de l'achèvement des attributs de l'ingénieur.
\end{itemize}

\stanceheading{Les recommandations aux partenaires, aux parties intéressées et aux autres entités}
\begin{itemize}
 \item La FCEG incite le BCAPG d'actualiser ses pratiques en matière de visites d'accréditation afin d'intégrer un retour d'information plus précieux et plus actif de la part des étudiants à la fin de chaque année d'accréditation.
 \item La FCEG incite le BCAPG à poursuivre ses efforts, notamment en ce qui concerne les Attributs des diplômés, afin de mieux refléter les résultats de l'apprentissage et le temps total d'apprentissage.
 \item La FCEG incite le BCAPG à réviser les Unités d'accréditation, car l'éducation et les concepts sont plus qu'un simple nombre quantifiable ou une case à cocher, car ils peuvent ne pas représenter avec précision la compréhension de l'étudiant.
 \item La FCEG incite le BCAPG à inclure les représentants des étudiants de la FCEG en tant que parties prenantes égales dans tous les projets, comités et groupes de travail liés à l'accréditation [réalisé].
 \item La FCEG incite les DDIC à collaborer sur des recommandations communes concernant le système d'accréditation actuel, basées sur les intérêts partagés des étudiants et des facultés.
 \item La FCEG incite les DDIC à avoir une communication transparente avec leurs écoles respectives afin que les voix des étudiants soient entendues au moment de la visite et que la meilleure formation d'ingénieur soit offerte.
 \item La FCEG incite le BCAPG à envisager d'intégrer davantage de possibilités d'expériences d'apprentissage extrascolaires dans le système d'accréditation actuel. Pour assurer la qualité de ces expériences d'apprentissage, la FCEG travaillera avec le BCAPG et les organismes provisoires de réglementation de l'ingénierie afin de créer un cadre et un ensemble de normes rigoureux.
 \item La FCEG incite le BCAPG de standardiser et de réguler la qualité des stages offerts par les programmes universitaires en ingénierie au Canada
 \item La FCEG incite les DDIC à régulièrement mettre à jour ses offres de cours et ses politiques pour les adapter à l'évolution des critères des Unités d'accréditation et aux améliorations apportées durant chaque année universitaire.
\end{itemize}

\stanceheading{Ce que fait la FCEG}
\begin{itemize}
 \item La FCEG assure la représentation des étudiants en tant qu'observateurs au sein du BCAPG depuis plusieurs années. En 2015, la responsabilité de plaidoyer auprès du BCAPG a été formalisée.
 \item La FCEF a assisté à une rencontre des DDIC pour la première fois en 2017, et a discuté d'une collaboration à venir sur des enjeux liés à l'accréditation.
 \item Le réunions du BCAPG accordent à la FCEG le droit de s'exprimer lors des sessions ouvertes, tel que décrit dans la \href{https://engineerscanada.ca/sites/default/files/2022-02/Board-Policy-Manual-Combined-e.pdf}{politique} ci-dessous [1]:
 \begin{itemize}
  \item 6.9.1 F.1 Observateurs présents lors des réunions (1) Le BCAPG doit inviter à ses réunions, à titre d'observateurs, les représentants suivants qui se verront accorder le droit de parole pendant les séances ouvertes : a) Le président ou la présidente de la Fédération Canadienne des Étudiant(e)s en génie (FCEG) ou son mandataire;
  \item De plus, la Politique 7.2 du Conseil d'Administration d'Ingénieurs Canada demande au BCAPG d'entretenir des relations avec la FCEG:
  \begin{itemize}
   \item Les étudiants en génie sont des parties prenantes importantes de l'agrément. En plus de solliciter la rétroaction des étudiants lors des visites d'évaluation de programmes, le Bureau d'agrément est tenu de maintenir des liens avec la FCEG et d'inviter un représentant à observer ses réunions et à lui soumettre un rapport. Tous les frais de déplacement de ce représentant sont assumés par Ingénieurs Canada.
  \end{itemize}
 \end{itemize}
 \item Le BCAPG a mis à jour son modèle de programme de visite d'accréditation basé sur les commentaires des parties prenantes, y compris la FCEG, en prévoyant plus de temps pour la participation des étudiants pendant les visites d'accréditation.
 \item La FCEG a inclus une section sur les programmes d'ingénierie dans l'enquête nationale de 2023 afin d'évaluer la connaissance générale du BCAPG ainsi que la compréhension et l'inclusion des Unités d'accréditation.
 \item Le BCAPG a été consulté pour les questions de l'enquête nationale de 2023 et pour la mise à jour de la documentation sur le cahier de positions.
\end{itemize}

\stanceheading{Ce que la FCEG envisage de faire}
\begin{itemize}
 \item La FCEG s'efforcera d'être incluse dans tous les projets et comités pertinents liés au BCAPG.
 \item La FCEG créera et distribuera en permanence une enquête nationale afin d'actualiser ses connaissances sur les expériences des étudiants en matière d'accréditation et de services fondés sur l'accréditation. Les questions de l'enquête seront consultées chaque année par le BCAPG afin d'éviter tout recoupement avec leur propre enquête.
 \item La FCEG actualisera sa position sur l'accréditation au fur et à mesure que des changements significatifs interviendront dans la politique et la réglementation.
 \item La FCEG demandera aux DDIC de rendre compte des changements non suivis au sein de leurs universités en ce qui concerne les nouvelles politiques et réglementations spécifiques aux attributs de l'ingénierie.
 \item La FCEG recueillera des commentaires supplémentaires spécifiquement liés aux visites d'accréditation afin d'examiner le point de vue des étudiants sur la qualité des visites d'accréditation qui ont eu lieu.
\end{itemize}

\newpage
